\paragraph{Selecting a feasible acquisition protocol.}
The reference implementation chooses among supplied certified protocols
using the query interface, scalar-reading and information-round limits, and
the requested priority. On the same 29 external pairwise cases at 1\%
internal error, it selects bulk in all 29 when range retrieval is admissible;
under a three-reading restricted interface it uses two adaptive readings in
22 cases and fixed-three in seven. A one-round restriction instead selects
fixed-three in all 29. Restricting both interfaces to two readings leaves
seven cases requiring escalation. The selector does not relax precision or
treat a missing catalogue certificate as a general impossibility result.

The operational wrapper binds each certificate to its model pair,
coefficients, meter, block, retained endpoints and source hashes. It rejects
readout envelopes for a different task, timestamp, feed or unit. Online
Resource~1, S8 retains all 116 protocol-selection decisions, callback tests,
exclusions and separate task-binding checks. These protect against accidental
mismatch; they do not authenticate a provider or calibrate a meter.

Reading charges and feedback delays favor different protocols. If a reserved
reading costs one unit and an information round costs $r\geq0$, the 29-case
fixed-three plan costs $87+29r$, while adaptive-where-certified with fixed
fallback costs $65+51r$. The latter is preferable for $r<1$ and fixed-three
for $r>1$. This exact catalogue frontier is a declared-cost sensitivity,
not a measured saving. S8 retains the full grid and all 306 independently
checked candidate evaluations. Actual public-API retrieval favors bulk:
the measured three-task mean I/O times are 0.331~s for bulk and 0.497~s for
adaptive retrieval. A shared-query objective does not change that interface
advantage.
